\documentclass[oneside,letterpaper,11pt]{book}
\usepackage[spanish]{babel}
\usepackage{graphicx}
\usepackage[margin=1in]{geometry}
\usepackage[svgnames]{xcolor}
\usepackage[sc]{mathpazo} 
\usepackage{newpxmath}
\usepackage{booktabs}
\usepackage{cite}
\usepackage{microtype}
\usepackage{hyperref}
\hypersetup{
    colorlinks=true,
    linkcolor=blue,
    filecolor=magenta,     
    urlcolor=cyan,
    }
\usepackage{fancyhdr}

\usepackage{subcaption}
 
\setlength{\headheight}{14.1pt}







\begin{document}

\renewcommand{\listtablename}{Índice de tablas}
\renewcommand{\tablename}{Tabla}
\renewcommand{\figurename}{Figura}


\frontmatter
    \pagestyle{plain}
    \begin{titlepage}
    \begin{center}
        \vspace*{0pt}
        \includegraphics[width=0.8\textwidth]{images/logoUSM.pdf}
        
        \Large
        Trabajo de Tesis\\
        \huge
        \vspace{0.5cm}
        
        \textbf{Design and Sizing of an Energy Storage System for a Hybrid Tugboat}\\
        \vspace{1.5cm}
        
        \normalsize
        Tesis para optar al título de\\
        \textbf{Magister en Ciencias de la Ingeniería Electrónico}
        
        \vfill
            
        Alumno\\
        \textbf{Leonardo Solis Zamora}
        \vspace{0.8cm}
            
        Profesor Guía\\
        \textbf{Dr. Marcelo Pérez Leiva}
        \vspace{0.8cm}
        
        Comisión evaluadora\\
        \textbf{Nombre del primer correferente}, Correferente, UTFSM\\
        \textbf{Nombre del segundo correferente}, Correferente, CODELCO
        \vspace{1cm}
        
        Enero 2025		, Valparaíso, Chile
    \end{center}
\end{titlepage}    
    \vspace*{5cm}
\begin{flushright}
    \emph{This is the dedicatory page.}
\end{flushright}
    
    
\chapter{Agradecimientos}
This is the abstract
    \chapter{Resumen}
This is the abstract
    \chapter{Abstract}
This is the abstract Holaaaa

    \tableofcontents
    \listoffigures
    \listoftables

\mainmatter
    \pagestyle{fancy}
    \fancyhf{}
    \renewcommand{\chaptermark}[1]{ \markboth{#1}{} }
    \lhead{\fancyplain{}{\chaptername} \fancyplain{}{\thechapter}. \fancyplain{}{\leftmark}}
    \rhead{ \fancyplain{}{\thepage} }
    \renewcommand{\headrulewidth}{1pt} 
    \renewcommand{\footrulewidth}{0pt} % sin línea
    \chapter{Introducción}
Introducción

\section{Motivation and Background}


\section{Challenges and Research Opportunities}


\section{Thesis Objectives and Outline}




\cite{10735355},
\cite{10679076},
\cite{10472698},
\cite{10382632},
\cite{10269751},
\cite{10091779},
\cite{10136711},
\cite{10015728},
\cite{9937167},
\cite{9738734},
\cite{9913469},
\cite{9360490},
\cite{9340547},
\cite{9242285},
\cite{9184122},
\cite{9299429},
\cite{9046048},
\cite{8322257},
\cite{7945527},
\cite{8013775},
\cite{7956250},
\cite{7836306},
\cite{7604136},




\backmatter
    \clearpage 
\addcontentsline{toc}{chapter}{\bibname}
\bibliographystyle{IEEEtran} 
\bibliography{content/references} 
     
\end{document}